\section{Block selection, reordering and duplication}
\label{app:block-ops}

\Cref{sec:weight_tied} trains a stage to use a single function.
Here we ask the opposite question on the trained untied ConvNeXt-T, without any
retraining: can a subset of its stage-3 blocks, reordered or reused, replace the nine
blocks?
Each configuration is a sequence of block indices applied with one Euler step
$\Delta t$, so the horizon is $T{=}L\,\Delta t$ for a sequence of length $L$
(\cref{tab:app-block-ops}; full ImageNet validation).
The same evaluation run gives $80.79$, $80.53$ and $80.36\%$ for the refined schedules
$R{=}2,4,10$, matching \cref{tab:r-ablation} within $0.01$ points, so it uses the
same checkpoint.
Script: \texttt{scripts/convnext\_interpolation.py}.

\begin{table}[t]
    \centering
    \caption{Untied ConvNeXt-T stage~3 replaced by sequences of its own blocks, top-1
    (\%). $T$: horizon $L\,\Delta t$. Native nine blocks: $80.96$.}
    \label{tab:app-block-ops}
    \small
    \setlength{\tabcolsep}{3pt}
    \begin{tabular}{lrrr}
        \toprule
        Blocks & $\Delta t$ & $T$ & top-1 \\
        \midrule
        best single block ($5$)           & $9$   & $9$    & 39.67 \\
        best single block ($5$)           & $1$   & $1$    & 43.12 \\
        $0,8$ / $8,0$                     & $4.5$ & $9$    & 49.03 / 15.31 \\
        $1,7$ / $7,1$                     & $4.5$ & $9$    & 23.14 / 15.12 \\
        $2,6$ / $6,2$                     & $4.5$ & $9$    & 13.81 / 3.97 \\
        $3,5$ / $5,3$                     & $4.5$ & $9$    & 57.57 / 35.18 \\
        $4,4$                             & $4.5$ & $9$    & 30.73 \\
        $0,4,8$                           & $4.5$ & $13.5$ & 38.16 \\
        $0,4,8$                           & $3$   & $9$    & 60.62 \\
        $0,4,8$                           & $2.25$& $6.75$ & 66.57 \\
        $(0,4,8)\times3$                  & $1$   & $9$    & 52.42 \\
        $0,0,2,2,4,4,6,6,8$               & $1$   & $9$    & 74.69 \\
        $0,0,0,3,3,3,6,6,6$               & $1$   & $9$    & 55.13 \\
        $0,0,0,0,4,8,8,8,8$               & $1$   & $9$    & 54.60 \\
        \bottomrule
    \end{tabular}
\end{table}

No single block suffices: the best one, block $5$, reaches $39.7\%$ when integrated
over the whole horizon in one step and $43.1\%$ with a unit step, and block $6$ peaks
at $29.2\%$ near $\Delta t{=}1$ over a sweep from $0.1$ to $20$.
With two blocks the original order wins in all four pairs we tested in both orders,
by $8.0$ to $33.7$ points.
Using every other block twice keeps $74.7\%$, the best substitute we found, whereas
cycling blocks $0,4,8$ three times gives $52.4\%$.
The blocks are therefore not interchangeable copies of one operation, consistent
with the nearly orthogonal increments of \cref{sec:coherence} and with
the finding that repeating one middle layer is the most harmful intervention in
pretrained transformers~\citep{sun2025painters}.
